\subsection{Where support must move when the bottom is exact}\label{sec:zero-bottom-tradeoff}
The unrestricted construction can remove supported zeros from the bottom, but it cannot remove all of the support enlargement. The next lower bound applies to every lift with exact bottom supports, whether or not it has an effective refinement.

\begin{proposition}\label{prop:zero-bottom-tradeoff}
Let an associative strict-unit algebra on the stated ten-element basis over $\mathcal S$ have the rank-three trace table. Suppose that each coefficient at the output $\Sigma$ has exactly its nonzero trace support. For every $w\in\{1,2,12,21\}$, put $h_w=|w|-1$. Each of the two ordered products
\[
 \gamma_w^1\Sigma,\qquad \Sigma\gamma_w^1
\]
has a supported zero-trace weight at some cell output of degree at least $3+h_w$, and another at some cell output of degree at most $-3-h_w$. The two cell outputs need not be the same. Consequently there are at least sixteen supported zero-trace weights outside the bottom output. If the algebra has an effective refinement, these weights account for at least sixteen cancelling pairs.
\end{proposition}

\begin{proof}
Write $K_{xy}^v\subset\mathbb Z$ for the finite occupied support of the coefficient of $v$ in $xy$. Addition of supports is union, and multiplication is Minkowski addition. Thus coefficientwise associativity is an equality between unions of such sums, and a weight contributed by one intermediate basis vector cannot be removed by another.

Fix $y=\gamma_w^1$, let $i$ be the first letter of $w$, and put $x=\gamma_i^1$. The matrix identity $P_iP_w=DP_w$ gives the trace coefficient
\[
 c_{xy}^y=\delta_1D=-D(D-1)(D-4).
\]
Its Laurent coefficients at both degrees $3$ and $-3$ are $-1$, so $\{-3,3\}\subseteq K_{xy}^y$. The coefficient of $\Sigma$ in $y\Sigma$ is
\[
 c_{y\Sigma}^{\Sigma}=4D^{1+h_w}(D-1).
\]
Every Laurent coefficient on $[-2-h_w,2+h_w]$ is positive. The exact-bottom hypothesis therefore gives
\[
 K_{y\Sigma}^{\Sigma}=[-2-h_w,2+h_w]\cap\mathbb Z,
 \qquad K_{x\Sigma}^{\Sigma}=[-2,2]\cap\mathbb Z.
\]
The intermediate output $y$ in $(xy)\Sigma$ consequently contributes both weights $5+h_w$ and $-5-h_w$ to the final output $\Sigma$.

Consider the same coefficient in $x(y\Sigma)$. The intermediate output $\Sigma$ contributes only weights in $[-4-h_w,4+h_w]$. The intermediate unit contributes no weight, since the strict-unit coefficient of $\Sigma$ in $x\cdot1=x$ is absent. For every cell basis vector $z\in\Gamma$, the trace table gives
\[
 K_{xz}^{\Sigma}\subseteq[-2,2]\cap\mathbb Z.
\]
Indeed $x$ is a length-one first-cell vector. Against a length-one vector, the bottom radius is at most one; against a length-two vector, it is at most two, as follows directly from the bottom factors $W_0,U_1,V_1$ in Lemma~\ref{lem:table}. Exactness of all bottom supports gives the displayed inclusion for the lifted table as well.

Associativity at the weight $5+h_w$ now requires a cell vector $z$ and weights
\[
 r\in K_{y\Sigma}^z,\qquad s\in K_{xz}^{\Sigma},
 \qquad r+s=5+h_w.
\]
Since $s\le2$, one has $r\ge3+h_w$. Applying the same argument at $-5-h_w$ gives a possibly different cell vector and a weight $r\le-3-h_w$ in its coefficient of $y\Sigma$. Every cell-output trace coefficient of $y\Sigma$ vanishes identically. Both occupied weights are therefore supported zeros.

For the reversed product, let $i$ be the last letter of $w$ and put $x=\gamma_i^1$. The identity $P_wP_i=DP_w$ gives $c_{yx}^y=\delta_1D$. Apply the preceding argument to
\[
 \Sigma(yx)=(\Sigma y)x.
\]
The same bottom-radius bound holds for $K_{zx}^{\Sigma}$, so $\Sigma y$ also has the two asserted supported zero weights.

The four words and the two multiplication orders give eight distinct ordered input pairs. In each pair the positive and negative weights just found are distinct. Counting a supported weight separately for each ordered input pair and output basis vector therefore gives at least sixteen slots. In an effective refinement a supported zero has equal, positive sign multiplicities, so each of those slots contains at least one cancelling pair.
\end{proof}

This proposition is a necessary condition, not an effective construction or an impossibility theorem for exact bottom supports. It leaves open whether the lower bound can be attained and whether any effective lift with exact bottom supports exists.
